% GENERATED FILE. Edit canonical lessons, then run npm run generate:books.
% canonical-source-hash: fnv1a64:10f6cfaaf22fe6fa
% canonical-lessons: TA-C124-takkali, TA-C124-milakay, TA-C124-inci, TA-C124-puntu, TA-C124-ten

\chapter{Tomato, Chilli, Ginger, Garlic, Honey}
\label{ch:ta-tomato-chilli-ginger-garlic-honey}
\hlchaptermodality{124}

\begin{chapteropening}
\textbf{By the end of this chapter:} \emph{I can say a tomato, a chilli, ginger, garlic and honey.}

Complete the last lesson of chapter 124: I can say a tomato, a chilli, ginger, garlic and honey.
\end{chapteropening}

\section[\texorpdfstring{\ta{தக்காளி} (takkāḷi)}{தக்காளி (takkāḷi)}]{\ta{தக்காளி} (takkāḷi) — a tomato}

\label{lesson:TA-C124-takkali}



\begin{quote}
Before the new one: say the Tamil for a coconut, then the Tamil for an onion.
\end{quote}

\subsection*{You'll want to know: \ta{தக்காளி}}
\textbf{\ta{தக்காளி}} — \emph{takkāḷi} — \textquotedblleft{}a tomato\textquotedblright{}.

\textbf{In the family, and next door}

\noindent
\begin{tabularx}{\linewidth}{@{}>{\raggedright\arraybackslash}X>{\raggedright\arraybackslash}X>{\raggedright\arraybackslash}X@{}}
\toprule
\textbf{} & \textbf{word} & \textbf{same root} \\
\midrule
Telugu & \te{టమాటా} (\emph{ṭamāṭā}) &  \\
Malayalam & \ml{തക്കാളി} (\emph{takkāḷi}) & yes \\
Hindi (neighbour) & \dv{टमाटर} (\emph{ṭamāṭar}) &  \\
English & a tomato &  \\
\bottomrule
\end{tabularx}

\begin{grammarlens}[title={What you've built}]
One more word for this chapter.
\end{grammarlens}

\subsection*{Your turn}
\begin{itemize}
\raggedright
  \item \emph{Say it:} \emph{takkāḷi}
  \item \emph{Say it:} \emph{takkāḷi}, once more
  \item \emph{Say it:} say \emph{takkāḷi}
  \item \emph{Recall:} say the Tamil for a coconut, then the Tamil for an onion, then say \emph{takkāḷi} again
\end{itemize}

\begin{tcolorbox}[breakable,colback=teal!4,colframe=teal!35!black,title={Before you move on}]
What does \ta{தக்காளி} mean? (\textquotedblleft{}A tomato\textquotedblright{}.) Say it once more.
\end{tcolorbox}
\section[\texorpdfstring{\ta{மிளகாய்} (milakāy)}{மிளகாய் (milakāy)}]{\ta{மிளகாய்} (milakāy) — a chilli}

\label{lesson:TA-C124-milakay}



\begin{quote}
Before the new one: say the Tamil for an onion, then the Tamil for a tomato.
\end{quote}

\subsection*{You'll want to know: \ta{மிளகாய்}}
\textbf{\ta{மிளகாய்}} — \emph{milakāy} — \textquotedblleft{}a chilli\textquotedblright{}.

\textbf{In the family, and next door}

\noindent
\begin{tabularx}{\linewidth}{@{}>{\raggedright\arraybackslash}X>{\raggedright\arraybackslash}X>{\raggedright\arraybackslash}X@{}}
\toprule
\textbf{} & \textbf{word} & \textbf{same root} \\
\midrule
Telugu & \te{మిరప} (\emph{mirapa}) & yes \\
Malayalam & \ml{മുളക്} (\emph{muḷakŭ}) & yes \\
Hindi (neighbour) & \dv{मिर्च} (\emph{mirc}) &  \\
English & a chilli &  \\
\bottomrule
\end{tabularx}

\begin{grammarlens}[title={What you've built}]
One more word for this chapter.
\end{grammarlens}

\subsection*{Your turn}
\begin{itemize}
\raggedright
  \item \emph{Say it:} \emph{milakāy}
  \item \emph{Say it:} \emph{milakāy}, once more
  \item \emph{Say it:} say \emph{milakāy}
  \item \emph{Recall:} say the Tamil for an onion, then the Tamil for a tomato, then say \emph{milakāy} again
\end{itemize}

\begin{tcolorbox}[breakable,colback=teal!4,colframe=teal!35!black,title={Before you move on}]
What does \ta{மிளகாய்} mean? (\textquotedblleft{}A chilli\textquotedblright{}.) Say it once more.
\end{tcolorbox}
\section[\texorpdfstring{\ta{இஞ்சி} (iñci)}{இஞ்சி (iñci)}]{\ta{இஞ்சி} (iñci) — ginger}

\label{lesson:TA-C124-inci}



\begin{quote}
Before the new one: say the Tamil for a tomato, then the Tamil for a chilli.
\end{quote}

\subsection*{You'll want to know: \ta{இஞ்சி}}
\textbf{\ta{இஞ்சி}} — \emph{iñci} — \textquotedblleft{}ginger\textquotedblright{}.

\textbf{In the family, and next door}

\noindent
\begin{tabularx}{\linewidth}{@{}>{\raggedright\arraybackslash}X>{\raggedright\arraybackslash}X>{\raggedright\arraybackslash}X@{}}
\toprule
\textbf{} & \textbf{word} & \textbf{same root} \\
\midrule
Kannada & \ka{ಶುಂಠಿ} (\emph{śuṇṭhi}) &  \\
Telugu & \te{అల్లం} (\emph{allaṁ}) &  \\
Malayalam & \ml{ഇഞ്ചി} (\emph{iñci}) & yes \\
Hindi (neighbour) & \dv{अदरक} (\emph{adrak}) &  \\
English & ginger &  \\
\bottomrule
\end{tabularx}

\begin{grammarlens}[title={What you've built}]
One more word for this chapter.
\end{grammarlens}

\subsection*{Your turn}
\begin{itemize}
\raggedright
  \item \emph{Say it:} \emph{iñci}
  \item \emph{Say it:} \emph{iñci}, once more
  \item \emph{Say it:} say \emph{iñci}
  \item \emph{Recall:} say the Tamil for a tomato, then the Tamil for a chilli, then say \emph{iñci} again
\end{itemize}

\begin{tcolorbox}[breakable,colback=teal!4,colframe=teal!35!black,title={Before you move on}]
What does \ta{இஞ்சி} mean? (\textquotedblleft{}Ginger\textquotedblright{}.) Say it once more.
\end{tcolorbox}
\section[\texorpdfstring{\ta{பூண்டு} (pūṇṭu)}{பூண்டு (pūṇṭu)}]{\ta{பூண்டு} (pūṇṭu) — garlic}

\label{lesson:TA-C124-puntu}



\begin{quote}
Before the new one: say the Tamil for a chilli, then the Tamil for ginger.
\end{quote}

\subsection*{You'll want to know: \ta{பூண்டு}}
\textbf{\ta{பூண்டு}} — \emph{pūṇṭu} — \textquotedblleft{}garlic\textquotedblright{}.

\textbf{In the family, and next door}

\noindent
\begin{tabularx}{\linewidth}{@{}>{\raggedright\arraybackslash}X>{\raggedright\arraybackslash}X@{}}
\toprule
\textbf{} & \textbf{word} \\
\midrule
Kannada & \ka{ಬೆಳ್ಳುಳ್ಳಿ} (\emph{beḷḷuḷḷi}) \\
Telugu & \te{వెల్లుల్లి} (\emph{vellulli}) \\
Malayalam & \ml{വെളുത്തുള്ളി} (\emph{veḷuttuḷḷi}) \\
Hindi (neighbour) & \dv{लहसुन} (\emph{lahsun}) \\
English & garlic \\
\bottomrule
\end{tabularx}

\begin{grammarlens}[title={What you've built}]
One more word for this chapter.
\end{grammarlens}

\subsection*{Your turn}
\begin{itemize}
\raggedright
  \item \emph{Say it:} \emph{pūṇṭu}
  \item \emph{Say it:} \emph{pūṇṭu}, once more
  \item \emph{Say it:} say \emph{pūṇṭu}
  \item \emph{Recall:} say the Tamil for a chilli, then the Tamil for ginger, then say \emph{pūṇṭu} again
\end{itemize}

\begin{tcolorbox}[breakable,colback=teal!4,colframe=teal!35!black,title={Before you move on}]
What does \ta{பூண்டு} mean? (\textquotedblleft{}Garlic\textquotedblright{}.) Say it once more.
\end{tcolorbox}
\section[\texorpdfstring{\ta{தேன்} (tēṉ)}{தேன் (tēṉ)}]{\ta{தேன்} (tēṉ) — honey}

\label{lesson:TA-C124-ten}



\begin{quote}
Before the new one: say the Tamil for ginger, then the Tamil for garlic.
\end{quote}

\subsection*{You'll want to know: \ta{தேன்}}
\textbf{\ta{தேன்}} — \emph{tēṉ} — \textquotedblleft{}honey\textquotedblright{}.

\textbf{In the family, and next door}

\noindent
\begin{tabularx}{\linewidth}{@{}>{\raggedright\arraybackslash}X>{\raggedright\arraybackslash}X>{\raggedright\arraybackslash}X@{}}
\toprule
\textbf{} & \textbf{word} & \textbf{same root} \\
\midrule
Kannada & \ka{ಜೇನು} (\emph{jēnu}) & yes \\
Telugu & \te{తేనె} (\emph{tēne}) & yes \\
Hindi (neighbour) & \dv{शहद} (\emph{śahad}) &  \\
English & honey &  \\
\bottomrule
\end{tabularx}

\begin{grammarlens}[title={What you've built}]
One more word for this chapter.
\end{grammarlens}

\subsection*{Your turn}
\begin{itemize}
\raggedright
  \item \emph{Say it:} \emph{tēṉ}
  \item \emph{Say it:} \emph{tēṉ}, once more
  \item \emph{Say it:} say \emph{tēṉ}
  \item \emph{Recall:} say the Tamil for ginger, then the Tamil for garlic, then say \emph{tēṉ} again
\end{itemize}

\begin{tcolorbox}[breakable,colback=teal!4,colframe=teal!35!black,title={Before you move on}]
What does \ta{தேன்} mean? (\textquotedblleft{}Honey\textquotedblright{}.) Say it once more.
\end{tcolorbox}
